% ============================================================
%  Math Notes Template
% ============================================================
\documentclass[11pt, a4paper]{article}

% --- Mathematics (must load before newtxmath) ---
\usepackage{amsmath, amssymb, amsthm, mathtools}

% --- Fonts & Encoding (XeLaTeX) ---
\usepackage{fontspec}
\usepackage{libertine}                    % Linux Libertine text + Biolinum sans
\usepackage[libertine]{newtxmath}         % Matching math font (loads after amsthm)
\usepackage{inconsolata}                  % Monospace font
\usepackage{microtype}                    % Improved typography

% --- Page Layout ---
\usepackage[
  margin=2.5cm,
  headheight=14pt
]{geometry}

% --- Line Spacing ---
\usepackage{setspace}
\onehalfspacing

% --- Paragraph Spacing ---
\usepackage{parskip}

% --- Colours ---
\usepackage[dvipsnames]{xcolor}
\definecolor{defblue}{HTML}{1a5276}
\definecolor{thmgreen}{HTML}{196f3d}
\definecolor{exred}{HTML}{922b21}
\definecolor{linkblue}{HTML}{2471a3}

% --- TikZ ---
\usepackage{tikz}
\usetikzlibrary{
  arrows.meta,
  calc,
  decorations.markings,
  patterns,
  positioning,
  shapes.geometric,
  cd                                      % commutative diagrams
}

% --- Theorem Environments ---
\usepackage{thmtools}

\declaretheoremstyle[
  headfont=\bfseries\color{thmgreen},
  bodyfont=\normalfont,
  spaceabove=10pt,
  spacebelow=10pt,
]{thmstyle}

\declaretheoremstyle[
  headfont=\bfseries\color{defblue},
  bodyfont=\normalfont,
  spaceabove=10pt,
  spacebelow=10pt,
]{defstyle}

\declaretheoremstyle[
  headfont=\bfseries\color{exred},
  bodyfont=\normalfont,
  spaceabove=10pt,
  spacebelow=10pt,
]{exstyle}

\declaretheorem[style=thmstyle, numberwithin=section, name=Theorem]{theorem}
\declaretheorem[style=thmstyle, sibling=theorem, name=Lemma]{lemma}
\declaretheorem[style=thmstyle, sibling=theorem, name=Proposition]{proposition}
\declaretheorem[style=thmstyle, sibling=theorem, name=Corollary]{corollary}
\declaretheorem[style=defstyle, sibling=theorem, name=Definition]{definition}
\declaretheorem[style=exstyle,  sibling=theorem, name=Example]{example}
\declaretheorem[style=exstyle,  numbered=no,      name=Remark]{remark}

% --- Hyperlinks & Cross-references ---
\usepackage[
  colorlinks=true,
  linkcolor=linkblue,
  citecolor=linkblue,
  urlcolor=linkblue
]{hyperref}
\usepackage{cleveref}

% --- Headers & Footers ---
\usepackage{fancyhdr}
\pagestyle{fancy}
\fancyhf{}
\fancyhead[L]{\textsc{\nouppercase{\leftmark}}}
\fancyhead[R]{\thepage}
\renewcommand{\headrulewidth}{0.4pt}

% --- Enumeration ---
\usepackage{enumitem}

% --- Bibliography (optional, uncomment if needed) ---
% \usepackage[style=alphabetic, backend=biber]{biblatex}
% \addbibresource{refs.bib}

% ============================================================
%  Fill these in for each set of notes
% ============================================================
\newcommand{\notestitle}{Hatcher's Algebraic Topology}
\newcommand{\notesauthor}{B. Arm}
\newcommand{\notesdate}{\today}

\title{\notestitle}
\author{\notesauthor}
\date{\notesdate}

% ============================================================
\begin{document}

\maketitle
\tableofcontents
\newpage

% ============================================================
%  Chapter 0: Some Underlying Geometric Notions
% ============================================================
\section{Some Underlying Geometric Notions}

\subsection{Homotopy and Homotopy Type}
\begin{definition}
  \textbf{Deformation Retraction: }A deformation retraction of a space $X$ onto a subspace $A\subset X$ is a family of maps $f_t: X\to X, t\in I$ such that $f_0 = \text{id}_X$, $f_1(X) = A$, $f_t|_A = \text{id}_A$.  
\end{definition}

\begin{definition}
  \textbf{Mapping Cylinder: }For a map $f: X\to Y$ the mapping cylinder $M_f$ is ther quotient space of the disjoint union $(X\times I)\sqcup Y$ obtained via identifying $X\times I \ni (x,1) \sim f(x) \in Y$:
  \[M_f:= \left((X\times I)\sqcup Y\right) / \sim\]
\end{definition}

A deformation retraction is a special case of the notion of a homotopy. 

\begin{definition}
  \textbf{Homotopy: }A homotopy is a family of maps $f_t: X\to Y, t\in I$ such that the associated map $F(x_t) = f_t(x)$ is continuous. Further one says to functions $f_0, f_1$ are \textit{homotopic} if there is a homotopy $f_t$ connecting them.
\end{definition}

\subsection{Cell Complexes}

\subsection{Operations on Spaces}

\subsection{Two Criteria for Homotopy Equivalence}

\subsection{The Homotopy Extension Property}

% ============================================================
%  Chapter 1: The Fundamental Group
% ============================================================
\section{The Fundamental Group}
The fundamental group of a space seeks to formalize the notion of probing for 'holes' in the space via the identification via homotopy equivalence of loops. 

% ------------------------------------------------------------
\subsection{Basic Constructions}
% ------------------------------------------------------------
\subsubsection{Paths and Homotopy}

\begin{definition}
  \textbf{Path: }A path in a space $X$ is a continious map $f: I\to X$.
\end{definition}

\begin{definition}
  \textbf{Path Homotopy: }A homotopy of paths in $X$ is a family $f_t: I\to X, t\in I$  such that:
  \begin{itemize}
    \item $f_t(0) = x_0, f_t(1) = x_1$ $\forall t$.
    \item The associated map $F(s,t) = f_t(s)$ is continuous. 
  \end{itemize}
\end{definition}

\begin{proposition}
  Homotopy equivalence on paths is an equivalence relation 
\end{proposition}
\proof 
\textbf{Reflexivity: }Clearly $f\simeq f$ via the constant homotopy $f_t = f$. 

\textbf{Symmetry: } Let $f\stackrel{\phi_t}{\simeq} g$, then $g\stackrel{\phi_{1-t}}{\simeq} f$ trivially.

\textbf{Transitivity: }Let $f\stackrel{\phi^{f,h}_t}{\simeq} g, g\stackrel{\phi^{g,h}_t}{\simeq} h$ we want to prove $f\simeq h$. We can do this via the homotopy $\phi^{f,h}_t = \mathbb{I}\{t\in [0,\frac12]\}\phi^{f,g}_{2t} + \mathbb{I}\{t\in (\frac12, 1]\}\phi^{g,h}_{2t-1}$.

A loop is a path $f: I\to X$ such that $f(0) = f(1) = x_0\in C$ and $x$ is referred to as the basepoint of the loop $f$. The set of homotopy classes of loops in $X$ with basepoint $x_0$ is denoted $\pi_1(X, x_0)$ and is named the fundamental group of the space $X$. 

For example if $X\subseteq \mathbb{R}^n$ is some convex subspace then trivially $\pi_1(X, x) = 0$ for any point $x\in X$. 

A more interesting point is that of the independence of the fundemental group from the chosen basepoint. In fact if the space is path connected then one can easily prove that $\pi_1(X, x_0) \cong \pi_1(X,x_1)$ for any points $x_0, x_1 \in X$.

\subsubsection{The Fundamental Group of the Circle}

\subsubsection{Induced Homomorphisms}

% ------------------------------------------------------------
\subsection{Van Kampen's Theorem}
% ------------------------------------------------------------

\subsubsection{Free Products of Groups}

\subsubsection{The Van Kampen Theorem}

\subsubsection{Applications to Cell Complexes}

% ------------------------------------------------------------
\subsection{Covering Spaces}
% ------------------------------------------------------------

\subsubsection{Lifting Properties}

\subsubsection{The Classification of Covering Spaces}

\subsubsection{Deck Transformations and Group Actions}

% ------------------------------------------------------------
\subsection*{Additional Topics}
% ------------------------------------------------------------

\subsubsection{Graphs and Free Groups}

\subsubsection{$K(G,1)$ Spaces and Graphs of Groups}

% ============================================================
%  Chapter 2: Homology
% ============================================================
\section{Homology}

% ------------------------------------------------------------
\subsection{Simplicial and Singular Homology}
% ------------------------------------------------------------

\subsubsection{$\Delta$-Complexes}

\subsubsection{Simplicial Homology}

\subsubsection{Singular Homology}

\subsubsection{Homotopy Invariance}

\subsubsection{Exact Sequences and Excision}

\subsubsection{The Equivalence of Simplicial and Singular Homology}

% ------------------------------------------------------------
\subsection{Computations and Applications}
% ------------------------------------------------------------

\subsubsection{Degree}

\subsubsection{Cellular Homology}

\subsubsection{Mayer--Vietoris Sequences}

\subsubsection{Homology with Coefficients}

% ------------------------------------------------------------
\subsection{The Formal Viewpoint}
% ------------------------------------------------------------

\subsubsection{Axioms for Homology}

\subsubsection{Categories and Functors}

% ------------------------------------------------------------
\subsection*{Additional Topics}
% ------------------------------------------------------------

\subsubsection{Homology and Fundamental Group}

\subsubsection{Classical Applications}

\subsubsection{Simplicial Approximation}

% ============================================================
%  Chapter 3: Cohomology
% ============================================================
\section{Cohomology}

% ------------------------------------------------------------
\subsection{Cohomology Groups}
% ------------------------------------------------------------

\subsubsection{The Universal Coefficient Theorem}

\subsubsection{Cohomology of Spaces}

% ------------------------------------------------------------
\subsection{Cup Product}
% ------------------------------------------------------------

\subsubsection{The Cohomology Ring}

\subsubsection{A K\"unneth Formula}

\subsubsection{Spaces with Polynomial Cohomology}

% ------------------------------------------------------------
\subsection{Poincar\'e Duality}
% ------------------------------------------------------------

\subsubsection{Orientations and Homology}

\subsubsection{The Duality Theorem}

\subsubsection{Connection with Cup Product}

\subsubsection{Other Forms of Duality}

% ------------------------------------------------------------
\subsection*{Additional Topics}
% ------------------------------------------------------------

\subsubsection{Universal Coefficients for Homology}

\subsubsection{The General K\"unneth Formula}

\subsubsection{$H$-Spaces and Hopf Algebras}

\subsubsection{The Cohomology of $SO(n)$}

\subsubsection{Bockstein Homomorphisms}

\subsubsection{Limits and Ext}

\subsubsection{Transfer Homomorphisms}

\subsubsection{Local Coefficients}

% ============================================================
%  Chapter 4: Homotopy Theory
% ============================================================
\section{Homotopy Theory}

% ------------------------------------------------------------
\subsection{Homotopy Groups}
% ------------------------------------------------------------

\subsubsection{Definitions and Basic Constructions}

\subsubsection{Whitehead's Theorem}

\subsubsection{Cellular Approximation}

\subsubsection{CW Approximation}

% ------------------------------------------------------------
\subsection{Elementary Methods of Calculation}
% ------------------------------------------------------------

\subsubsection{Excision for Homotopy Groups}

\subsubsection{The Hurewicz Theorem}

\subsubsection{Fiber Bundles}

\subsubsection{Stable Homotopy Groups}

% ------------------------------------------------------------
\subsection{Connections with Cohomology}
% ------------------------------------------------------------

\subsubsection{The Homotopy Construction of Cohomology}

\subsubsection{Fibrations}

\subsubsection{Postnikov Towers}

\subsubsection{Obstruction Theory}

% ------------------------------------------------------------
\subsection*{Additional Topics}
% ------------------------------------------------------------

\subsubsection{Basepoints and Homotopy}

\subsubsection{The Hopf Invariant}

\subsubsection{Minimal Cell Structures}

\subsubsection{Cohomology of Fiber Bundles}

\subsubsection{The Brown Representability Theorem}

\subsubsection{Spectra and Homology Theories}

\subsubsection{Gluing Constructions}

\subsubsection{Eckmann--Hilton Duality}

\subsubsection{Stable Splittings of Spaces}

\subsubsection{The Loopspace of a Suspension}

\subsubsection{The Dold--Thom Theorem}

\subsubsection{Steenrod Squares and Powers}

% ============================================================
%  Appendix
% ============================================================
\section*{Appendix}
\addcontentsline{toc}{section}{Appendix}

\subsection*{Topology of Cell Complexes}

\subsection*{The Compact-Open Topology}

\subsection*{The Homotopy Extension Property}

\subsection*{Simplicial CW Structures}

% \printbibliography

\end{document}
